%% kapa: exam=1 round=37 year=2026 period=2 form=A subject=appraisal_law subject_ko=감정평가관계법규 seq=17 number=17
%% source: Q-Net 감정평가사 기출문제 (cst00309) artlSeq=5258924 file=raw/1008/5258924/2026년 제37회 감정평가사 1차 2교시.pdf
%% conversion: kordoc 4.24.51 -> tools/gen1.py (auto); review=proofread-sonnet; flags=none
\begin{question}{17}
감정평가 및 감정평가사에 관한 법령상 감정평가법인등의 권리와 의무에 관한 설명으로 옳지 않은 것은?
\begin{choices}
  \item 감정평가사가 합동사무소를 설치하려는 경우 2명 이상의 감정평가사를 두어야 한다.
  \item 감정평가법인등의 업무수행에 관한 수수료의 요율 및 실비의 범위는 국토교통부장관이 감정평가관리ㆍ징계위원회의 심의를 거쳐 결정한다.
  \item 감정평가법인등은 손해배상책임의 보장을 위한 보증보험금으로 손해배상을 하였을 때에는 10일 이내에 보험계약을 다시 체결해야 한다.
  \item 금고 이상의 형의 집행유예를 받고 그 유예기간이 만료된 날부터 2년이 지난 감정평가사는 감정평가사사무소 개설의 결격사유에 해당한다.
  \item 감정평가법인 소속 사무직원은 소속된 감정평가법인 이외의 다른 감정평가법인의 주식을 소유할 수 없다.
\end{choices}
\end{question}
\answer{4}
